\section*{Заключение}
\addcontentsline{toc}{section}{Заключение}

В ходе выполнения четырёх лабораторных работ были реализованы все
поставленные задания. В лабораторной работе~1 выполнена случайная генерация
связного ациклического ориентированного графа по распределению Вейбулла,
вычислены экцентриситеты вершин, выделены центр и диаметральные вершины,
реализован метод Шимбелла для поиска кратчайших и наидлиннейших маршрутов
заданной длины, а также определение всех простых путей между двумя вершинами.
В лабораторной работе~2 реализованы топологическая сортировка по алгоритму
Фалкерсона, нахождение кратчайших путей алгоритмом Флойда-Уоршалла
с выводом матрицы расстояний и сравнение алгоритмов по числу итераций.
В лабораторной работе~3 построена транспортная сеть, найден максимальный поток
алгоритмом Форда--Фалкерсона и вычислен поток минимальной стоимости
величиной $\lfloor 2/3 \cdot F_{\max} \rfloor$.
В лабораторной работе~4 подсчитано число остовных деревьев по теореме Кирхгофа,
построен минимальный остов алгоритмом Краскала, выполнено кодирование
и декодирование кода Прюфера с сохранением весов рёбер, а также найдено
максимальное независимое множество вершин.

\subsection*{Достоинства реализации}

Все четыре лабораторные работы используют один и тот же сгенерированный граф
без перегенерации, что удобно при тестировании разных алгоритмов на одних
и тех же данных.

Алгоритм поиска кратчайшего пути по стоимости в третьей лабораторной
переиспользует тот же алгоритм Флойда-Уоршалла, что реализован во второй.
Это позволило не писать один и тот же код дважды.

Генерация весов рёбер и перемешивание вершин при построении графа
используют одно и то же распределение Вейбулла, что делает результат
более случайным и непредсказуемым по сравнению с равномерным генератором.

\subsection*{Недостатки реализации}

Граф хранится в виде матриц, поэтому при большом числе вершин программа
занимает много памяти, даже если рёбер мало.

В алгоритме Краскала проверка на цикл выполняется обходом уже построенного
остова, что работает медленнее, чем специальные структуры данных для
этой задачи.

Поиск всех путей и максимального независимого множества работает очень
долго на больших графах, так как перебирает все возможные варианты.
Никакого ограничения на размер входных данных не предусмотрено.

\subsection*{Масштабирование (что можно добавить в дальнейшем)}

Можно заменить матричное представление графа на списки смежности,
чтобы программа работала быстрее и занимала меньше памяти на разреженных
графах.

Добавление графического интерфейса позволило бы отображать граф наглядно:
рисовать вершины и рёбра, выделять цветом центр и диаметральные вершины,
показывать пошагово работу алгоритмов и менять параметры без перезапуска
программы.
